\documentclass[10pt,a4paper]{amsart}
\usepackage[margin=19mm]{geometry}
\usepackage{amsmath,amssymb,mathtools}
\usepackage[scr=rsfs,cal=euler]{mathalfa}
\usepackage{xfrac}
\usepackage[unicode,hidelinks,bookmarks=false]{hyperref}
\newcommand{\ie}{i.e.\ }
\newcommand{\cf}{cf.\ }
\newcommand{\resp}{respectively\ }
\newcommand{\Subsection}{\subsection}
\providecommand{\nobreakdash}{\nobreak\hbox{-}}
\setlength{\emergencystretch}{2em}
\multlinegap0pt
\usepackage{bm}
\usepackage{bbm}
\usepackage{dsfont}
\usepackage{xspace}
\usepackage{cancel}
\usepackage{mathtools}
\usepackage{amsthm}
\makeatletter
\@ifundefined{thm}{\newtheorem{thm}{Thm}}{}
\@ifundefined{prop}{\newtheorem{prop}[thm]{Proposition}}{}
\makeatother
\title[Homology Generators and Relations for the Ordered Configurat]{Homology Generators and Relations for the Ordered Configuration Space of a Star Graph}
\author{Nicholas Wawrykow}
\date{Source: arXiv:2401.13821v3 (CC BY 4.0)}
\begin{document}
\maketitle

\noindent\emph{Source and licence.} Homology Generators and Relations for the Ordered Configuration Space of a Star Graph. Version arXiv:2401.13821v3 (CC BY 4.0), distributed under the Creative Commons Attribution 4.0 International licence, as recorded in the arXiv licence metadata for this exact version and shown on the abstract page https://arxiv.org/abs/2401.13821v3. Full rights notice: licenses/arxiv-2401.13821-CC-BY-4.0.txt.\par\smallskip

\noindent\textit{Editorial scope.} Counted unit: the Prop whose source label is \texttt{conditions for FIOk space} in the article (source0.tex lines 300--303, source label \texttt{conditions for FIOk space}), delivered together with the article's own complete original proof of it (source0.tex lines 305--373, one \texttt{proof} environment of 69 lines). No mathematics is written, repaired or re-derived: statement and proof are the article's own text line for line. Required context retained uncounted: none. Every definition, display and label of those lines is retained unchanged. Omitted from this package and not redistributed: 1--299, 304--304, 374--1276. Nothing inside a retained excerpt is omitted or altered: every retained body is delivered exactly as the article's lines are written, so no omission receipt of any kind accompanies this package. Citations: the retained excerpts carry no citation command at all, so no bibliography entry of the article is transcribed and no reference is redistributed. Reference adaptation: none was needed and none was made; every reference inside the delivered unit resolves inside it, so no unresolved reference or citation is produced. Printed slips kept literal: no printed word, symbol, hypothesis or number of the counted statement or of its proof was altered. Layout and class: the article's own class is \texttt{article} with packages including geometry, tikz-cd, float, amsmath, amsfonts, amsthm, enumitem, dsfont, amssymb, pifont, mathtools, comment, graphicx, caption; the wrapper is the lane's pinned \texttt{amsart} editorial wrapper at 10pt on A4, which restates the theorem-like environments the retained text needs and adds the article's own macro definitions that the retained text needs (none), with extra wrapper packages (bm, bbm, dsfont, xspace, cancel, mathtools). The delivered numbering is the unit's own. Assets: no retained span contains a figure environment, a graphics-inclusion command, a PDF-page inclusion, a table or a TikZ picture; no image file is copied into this package, the package declares no asset dependency and no image file is redistributed with it. Licence: version arXiv:2401.13821v3 of this article is distributed under the Creative Commons Attribution 4.0 International licence, as recorded in the arXiv licence metadata for this exact version and shown on the abstract page https://arxiv.org/abs/2401.13821v3. A full rights notice is delivered at licenses/arxiv-2401.13821-CC-BY-4.0.txt. Characters: every retained excerpt is pure ASCII, so the delivered engine, XeLaTeX with its default Latin Modern text font, reports no missing character.
\begin{prop}\label{conditions for FIOk space}
Let $X$ be an FB-space such that for all $n$ there exist $d$ insertion maps $\iota_{n, j}:X_{n}\to X_{n+1}$.
If these insertion maps commute with permutations and distinct insertions are unordered, then $X$, along with these insertion maps, has the structure of an FI$_{d,o}$-space.
\end{prop}

\begin{proof}
For each morphism $(f, c, o):[n]\to [m]$ in FI$_{d,o}$, we need to define a map $(f, c, o)_{*}:X_{n}\to X_{m}$.
When $n=m$, we have $f=\sigma$ for some $\sigma\in S_{n}$, and $c=\cdot$ and $o=\cdot$ are the trivial coloring and ordering, respectively.
In this case $(f, c, o)_{*}=\sigma$, where we also use $\sigma$ to denote its action on $X_{n}$.
The insertion maps $\iota_{l,j_{l+1}}$ describe what happens when $f$ is the standard inclusion $\iota_{l}:[l]\hookrightarrow [l+1]$, $c_{j_{l+1}}$ colors the element $l+1$ color $j_{l+1}$, and $o_{j_{l+1}}$ orders the element $l+1$ the first among the elements colored $j=j_{l+1}$, i.e., $(f, c, o)=(i_{l}, c_{j_{l+1}}, o_{j_{l+1}})$ and $(f, c, o)_{*}=\iota_{l, j_{l+1}}$.

Given an arbitrary morphism $(f, c, o):[n]\to[m]$ in FI$_{d,o}$ we can write $f:[n]\to [m]$ as $\sigma'\circ \iota_{m-1, j_{m}}\circ\cdots\circ \iota_{n, j_{n+1}}$ for $\sigma'\in S_{m}$.
Moreover, there is a unique such $\sigma'$ that sends $1$ to $f(1)$, $2$ to $f(2)$, etc., while sending $n+1$ to the first element in complement of the image of $f$ colored $j_{n+1}$, $n+2$ to the first element colored $j_{n+2}$ complement of the image $f$ if $j_{n+2}\neq j_{n+1}$ and the second element colored $j_{n+1}$ if $j_{n+2}=j_{n+1}$, and so on and so forth.
It follows that 
\[
(f, c, o)=(\sigma', \cdot, \cdot)\circ(\iota_{m-1}, c_{j_{m}}, o_{j_{m}})\circ\cdots\circ (\iota_{n}, c_{j_{n+1}}, o_{j_{n+1}}),
\]
so we should define
\[
(f, c, o)_{*}:=\sigma'\circ\iota_{m-1, j_{m}}\circ\cdots\circ\iota_{n, j_{n+1}}.
\]

To check functoriality, we must show that if we have another sequence of permutations and insertions that compose to $(f, c, o)$ in FI$_{d,o}$, then the corresponding maps on the various spaces $X_{l}$ compose to $(f, c, o)_{*}$.
Given an arbitrary sequence of permutations and insertion maps, the property that insertion maps commute with permutations implies that we can push all the permutations to the left of the insertion maps without changing the composition map, obtaining a composition of insertions followed by a composition of permutations.
Since there is an $S_{n}$-action on $X_{n}$, we can replace the composition of permutations by a single permutation.
Therefore, it suffices to show that if
\[
(f, c, o)=(\sigma'', \cdot, \cdot)\circ (\iota_{m-1}, c_{j'_{m}}, o_{j'_{m}})\circ \cdots\circ(\iota_{n}, c_{j'_{n+1}}, o_{j'_{n+1}}),
\]
then
\[
\sigma''\circ\iota_{m-1, j'_{m}}\circ \cdots\circ\iota_{n, j'_{n+1}}=\sigma'\circ\iota_{m-1, j_{m}}\circ \cdots\circ\iota_{n, j_{n+1}}.
\]

We proceed by induction on $m-n$, noting that the above proves that this is true if $m-n=0$ or $1$.
Note that $\sigma'$ and $\sigma''$ take the same values on $[n]$ as $f$, so we can write $\sigma''=\sigma'\circ \tilde{\sigma}$ where $\tilde{\sigma}$ only permutes $[m]\backslash[n]$.
Therefore, after canceling $\sigma'$ from both sides, it suffices to show that 
\[
\tilde{\sigma}\circ\iota_{m-1, j'_{m}}\circ\cdots\circ\iota_{n, j'_{n+1}}=\iota_{m-1, j_{m}}\circ \cdots\circ\iota_{n, j_{n+1}}.
\]
This follows from the fact that distinct insertions are unordered.
We see this by induction on $m-n$.
If $m-n=1$, this is trivial.
Otherwise, let $n+k=(\tilde{\sigma})^{-1}(n+1)$, i.e., in the alternate composition, element $n+k$ gets inserted with color $j'_{n+k}=j_{n+1}$ and then $\tilde{\sigma}$ changes it to $n+1$.
Since distinct insertions are unordered if $j'_{n+k}\neq j'_{n+k-1}$, we can write
\[
\iota_{n+k-1, j'_{n+k}}\circ\iota_{n+k-2, j'_{n+k-1}}=(n+k\;\;\;n+k-1)\circ \iota_{n+k-1, j'_{n+k-1}}\circ\iota_{n+k-2, j'_{n+k}}.
\]
Moreover, we need not worry about the case $j'_{n+k}=j'_{n+k-1}$ as this would imply that the two compositions have a different first element of color $j'_{n+k}$ in the ordering given by $o$, a contradiction.
By applying the fact that insertion maps commute with permutations, we can move all the transpositions to the left to make the composition
\[
\tilde{\sigma}\circ (n+k\;\;\; n+k-1)\circ\cdots\circ (n+2\;\;\; n+1).
\]
This is equal to 
\[
\tilde{\sigma}\circ(n+k\;\;\;n+k-1\;\;\;\cdots\;\;\; n+2\;\;\; n+1),
\]
which fixes $n+1$.
Denote this composition by $\tilde{\sigma}'$.
It follows that
\[
\tilde{\sigma}\circ\iota_{m-1, j'_{m}}\circ\cdots\circ\iota_{n, j'_{n+1}}=\tilde{\sigma}'\circ \iota_{m-1, j'_{m}}\circ\cdots\circ\iota_{n+k, j'_{n+k+1}}\circ \iota_{n+k-1, j'_{n+k-1}}\circ\cdots\circ\iota_{n+1, j'_{n+1}}\circ\iota_{n, j'_{n+k}},
\]
and $j'_{n+k}=j_{n+1}$.
By the induction hypothesis, we have that 
\[
\tilde{\sigma}'\circ \iota_{m-1, j'_{m}}\circ\cdots\circ\iota_{n+k, j'_{n+k+1}}\circ\iota_{n+k-1, j'_{n+k-1}}\circ\cdots\circ\iota_{n+1, j'_{n+1}}
\]
and
\[
\iota_{m-1, j_{m}}\circ\cdots\circ\iota_{n+1, j_{n+2}}
\]
induce the same maps from $[n+1]$ to $[m]$ as they are both induced by $(f, c, o)$---here $n+1$ is not colored or ordered, rather it is sent to the element that is the first colored $j_{n+1}$ by  $(f, c, o)$---so by composing both of these on the right with $\iota_{n, j'_{n+k}}=\iota_{n, j_{n+1}}$ gives the desired equality, completing the inductive step and the proof of the proposition.
\end{proof}

\end{document}
